\section{Agent-native Toolchain}

ASIEP is implemented as an agent-native toolchain rather than only as a static schema. The goal is not merely to help a human auditor read a record, but also to let another agent inspect the record, classify errors, request missing evidence, generate bounded repair plans, resolve artifact references, package evidence, and rerun evaluation. This design reflects the expected operating environment of agent-to-agent systems, where validation and evidence preparation may be delegated to specialized agents. For this reason, the toolchain emphasizes stable machine-readable outputs rather than informal terminal messages.

The \texttt{asiep\_validator} is the primary conformance component. It checks JSON Schema conformance, required field groups, lifecycle state consistency, required evidence references, gate consistency, rollback evidence, and selected semantic invariants. Its output is structured JSON. Each error contains a stable error code, JSON path, invariant identifier when applicable, repairability class, and remediation hint. This allows a downstream agent to distinguish missing external evidence from unsafe promotion, missing rollback evidence, unresolved references, or digest mismatch. Stable error objects also make validator behavior reproducible across fixture versions.

The \texttt{asiep\_repairer} converts validation failures into bounded repair plans. It does not directly mutate the evidence object and does not fabricate missing evidence. This boundary is central to the design. A repair plan may recommend adding a missing reference, correcting a lifecycle state when supporting evidence already exists, or changing an invalid promotion decision to rejection. However, it must not invent a gate report, change a safety regression from true to false, forge a digest, or rewrite an artifact reference to hide tampering. Repairability is therefore treated as a constrained workflow: structural omissions may be repairable, but evidential failures require new evidence or human review.

The \texttt{asiep\_resolver} turns symbolic evidence references into locally verifiable bindings. It checks whether referenced artifacts exist inside a local evidence bundle, prevents path escape, recomputes SHA-256 digests, and reports missing artifacts or digest mismatches. This step prevents ASIEP from becoming a set of unverifiable links. A record that names a trace, feedback file, candidate diff, or gate report is not locally closed until the resolver confirms that the artifact exists, remains inside the bundle boundary, and matches the declared digest.

The \texttt{asiep\_importer} converts local OpenTelemetry-like and LangSmith-like fixtures into ASIEP evidence bundles. The importer follows a reference-only policy: it maps trace roots, spans, run identifiers, feedback records, scores, and candidate artifacts into ASIEP references rather than embedding all raw content into the profile object. This keeps the ASIEP record compact and separates the profile from artifact storage. The current importer is local and fixture-based. It does not claim full OpenTelemetry collector integration or real LangSmith API integration.

The \texttt{asiep\_packager} transforms validated and resolved bundles into local research-object outputs. These include a package manifest, an FDO-like local record, RO-Crate-like metadata, and PROV JSON-LD. The packager is intentionally conservative: it requires successful ASIEP validation and successful bundle resolution before packaging. Its outputs support local inspection, reproducibility, and cross-standard mapping, but they do not imply real FDO registry submission, externally resolvable identifier issuance, or full RO-Crate certification.

The \texttt{asiep\_evaluator} runs the local corpus and produces conformance and robustness metrics. It evaluates known valid records, invalid records, tampered bundles, missing evidence cases, gate violations, rollback cases, and cross-standard mapping coverage. It also generates paper-ready tables and reports. Its purpose is to test whether the current implementation behaves consistently on the local fixture corpus. It is not a general benchmark of agent capability and does not provide production assurance.

Together, these components make ASIEP more than a descriptive metadata template. The validator checks conformance, the repairer produces bounded repair plans, the resolver verifies local artifact bindings, the importer connects trace fixtures to evidence bundles, the packager creates local evidence packages, and the evaluator reports reproducible local results. The result is a runnable evidence-checking workflow for agent self-improvement events.
